\section{Návrh riešenia}

Návrh riešenia opisuje celkové usporiadanie systému a vzájomné prepojenie jeho hlavných častí. 
Systém je navrhnutý ako IoT riešenie, ktoré prepája fyzický laboratórny model, 
mikrokontrolér ESP32, lokálny Flask server, webové rozhranie, dátové úložisko a 
cloudovú platformu ThingsBoard.

\subsection{Celková architektúra systému}

Systém je rozdelený na štyri základné vrstvy. Fyzická vrstva obsahuje hydraulický okruh, 
snímač teploty a akčné členy, teda Peltierov článok, vodnú pumpu, ventilátor a vyhrievacie 
teleso. Riadiacu vrstvu tvorí mikrokontrolér ESP32, ktorý zabezpečuje čítanie teploty a 
generovanie PWM signálov pre výkonové prvky.

Lokálnu serverovú vrstvu tvorí počítač so spustenou Flask aplikáciou. Táto vrstva prijíma 
údaje z ESP32, spracúva aktuálny stav systému, poskytuje webové rozhranie a zabezpečuje 
ukladanie meraní. Cloudovú vrstvu tvorí platforma ThingsBoard, ktorá slúži na vzdialené 
sledovanie a vybrané ovládanie systému.

Takéto rozdelenie umožňuje oddeliť fyzickú realizáciu, lokálne riadenie, vizualizáciu a 
cloudový monitoring. Flask server pritom slúži ako centrálna komunikačná brána medzi ESP32, 
webovým rozhraním, databázou a ThingsBoardom.

\subsection{Bloková schéma systému}

Celkové usporiadanie systému je znázornené na blokovej schéme. Schéma zobrazuje hlavné 
funkčné vrstvy systému a spôsob ich vzájomného prepojenia.

\begin{figure}[H]
    \centering
    \safeimage{obrazky/Iot_architektura.png}{0.95\textwidth}{Bloková schéma IoT 
    architektúry systému riadenia teploty kvapaliny}
    \label{fig:blokova_schema_iot}
\end{figure}

Medzi fyzickou vrstvou a mikrokontrolérom ESP32 sa prenášajú merané údaje zo snímača a 
riadiace PWM signály pre výkonové členy. ESP32 komunikuje s lokálnym Flask serverom pomocou 
dátových správ. Flask server následne poskytuje údaje webovému rozhraniu, ukladá ich do 
databázy a JSON súboru a odosiela vybrané telemetrické údaje do platformy ThingsBoard.

\subsection{Tok dát a riadiacich príkazov}

Tok dát v systéme je navrhnutý tak, aby všetky merané hodnoty aj riadiace príkazy prechádzali 
cez Flask server. ESP32 teda priamo nekomunikuje s databázou ani s cloudovou platformou, 
ale odosiela dáta na lokálny server, ktorý ich ďalej spracúva.

\begin{center}
Snímač teploty $\rightarrow$ ESP32 $\rightarrow$ Flask server $\rightarrow$ webové rozhranie / databáza / ThingsBoard
\end{center}

Riadiace príkazy môžu prichádzať z webového rozhrania alebo z ThingsBoard dashboardu. Flask server ich uloží do aktuálneho stavu systému a pri ďalšej komunikácii ich odošle do ESP32.

\begin{center}
Webové rozhranie / ThingsBoard $\rightarrow$ Flask server $\rightarrow$ ESP32 $\rightarrow$ PWM výstupy
\end{center}

Podrobnejší opis konkrétnej komunikácie, endpointov, ukladania dát a cloudovej integrácie je 
uvedený v softvérovej časti dokumentácie.